Catalan 数 $C_n$ 出现在许多组合计数问题中，以下是几个经典例子：

\begin{enumerate}
    \item \textbf{路径计数问题}：有一个大小为 $n \times n$ 的方格图，左下角为 $(0,0)$，右上角为 $(n,n)$。从左下角出发，每次只能向右或向上走一单位，且不走到对角线 $y=x$ 上方（但可以触碰）的情况下，到达右上角的路径总数为 $C_n$。
    \item \textbf{圆内不相交弦计数问题}：圆上有 $2n$ 个点，将这些点成对连接起来且使得所得到的 $n$ 条线段两两不交的方案数为 $C_n$。
    \item \textbf{三角剖分计数问题}：对角线不相交的情况下，将一个凸 $(n+2)$ 边形区域分成三角形区域的方法数为 $C_n$。
    \item \textbf{二叉树计数问题}：含有 $n$ 个结点的形态不同的二叉树数目为 $C_n$；等价地，含有 $n$ 个非叶结点的形态不同的满二叉树数目也为 $C_n$。
    \item \textbf{括号序列计数问题}：由 $n$ 对括号构成的合法括号序列数为 $C_n$。
    \item \textbf{出栈序列计数问题}：一个栈（无穷大）的进栈序列为 $1,2,\dots,n$，合法出栈序列的数目为 $C_n$。
    \item \textbf{数列计数问题}：由 $n$ 个 $1$ 和 $n$ 个 $-1$ 组成的数列 $a_1,a_2,\dots,a_{2n}$，满足部分和 $a_1+a_2+\cdots+a_k \ge 0$ 对 $k=1,2,\dots,2n$ 均成立的数列数目为 $C_n$。
\end{enumerate}

Catalan 数有如下常见的表达式：
\[
C_n = \frac{1}{n+1}\binom{2n}{n} = \frac{(2n)!}{n!(n+1)!},\qquad n\ge 0 \tag{1}
\]
\[
C_n = \binom{2n}{n} - \binom{2n}{n+1},\qquad n\ge 0 \tag{2}
\]
\[
C_n = \frac{4n-2}{n+1}C_{n-1},\quad n>0,\quad C_0 = 1 \tag{3}
\]